\begin{example}[Cycle currents on the Berkovich line]
\label{ex:ch14-line-cycles}
Return to $X=\mathbf P^{1,\An}_k$ with $k$ trivially valued and
algebraically closed, as in \Cref{sec:ch14-line-examples}.  The analytic
cycles of dimension $0$ on $X$ are the locally finite formal sums
$Z=\sum_bm_b\,b$ over rigid points, and
\[
 \delta_Z=\sum_bm_b\,\delta_{x_b}
\]
is the corresponding measure; the sum is finite against any test function by
\Cref{lem:ch14-star-smooth}.  The only cycle of dimension $1$ is
$m\cdot X$, with $\delta_{mX}=m\,\delta_X$; by
\Cref{cor:ch14-delta-x-closed} it is closed, and its ``mass'' against the
constant test form is the total length of the skeleton of $X$, which is zero
because the skeleton of $\mathbf P^{1,\An}_k$ over a trivially valued field is
a point.  This is the degenerate case in which the calibration carries no
metric information; contrast \Cref{ex:ch14-skeleton-current}.

Finally, \Cref{ex:ch14-line-model-function} exhibits a difference of two cycle
currents as a $\dprimeddprime$: with $r=1$ and $c_1=1$, the identity there
reads
\[
 \dprimeddprime[u]=\delta_{p_1}-\delta_\eta,
\]
with $u=\log\max\{|T-b_1|,\rme^{-1}\}$ on the branch $b_1$ and $u=0$
elsewhere.  This is the truncated local model of the Poincar\'e--Lelong
formula \eqref{eq:appJ-cld-pl}; the untruncated statement, and the passage
from a truncation to a divisor, are the subject of
\Cref{chap:cld-poincare-lelong}.
\end{example}

\begin{theorem}
\label{thm:current-fundamental-cycle}
For every pure-dimensional closed analytic subspace $Y\subset X$,
\[
                         \delta_Y=\delta_{[Y]}.
\]
Thus nilpotent and embedded structure enters precisely through the generic
multiplicities of the fundamental cycle.
\end{theorem}
\begin{proof}
Write $Y_1,\ldots,Y_s$ for the irreducible components of
$Y_{\mathrm{red}}$ and let
$m_j=\operatorname{length}_{\mathcal O_{Y,\eta_j}}
\mathcal O_{Y,\eta_j}$ be the generic multiplicity along $Y_j$.  Remove the
union $B$ of the proper intersections $Y_i\cap Y_j$ and the loci where an
irreducible component is not generically reduced.  This is a closed analytic
set of dimension less than $r$.  A test form of bidegree $(r,r)$ restricts
to zero on $B$ by dimension, so deleting $B$ does not change either
integration current.

On $Y\setminus B$ the components are disjoint.  Choose a tropical chart at a
point of $Y_j\setminus B$.  The canonical calibration of the nonreduced
analytic space is defined by the initial cycle of its affinoid algebra.
Localization at the generic prime of $Y_j$ shows that this initial cycle is
$m_j$ times the initial cycle of $(Y_j)_{\mathrm{red}}$: multiplicity is the
length of the localized algebra, and length is unchanged by the graded
reduction used to define the calibration.  Hence chartwise integration over
$Y$ equals $m_j$ times integration over $Y_j$.

Summing over the disjoint components gives
\[
          \delta_Y=\sum_jm_j\delta_{Y_j}=\delta_{[Y]}.
\]
The removal and graded-length arguments are
\cite[Proposition~16.6.3]{CLD25}.  They also show why replacing $Y$ by its
reduction without the coefficients $m_j$ would be incorrect.
\end{proof}
